%
% 2-wurzel.tex
%
% (c) 2026 Prof Dr Andreas Müller
%
\section{Integrale von $R(x,\vartheta)$ mit $\vartheta=\!\sqrt{ax^2+bx+c}$
\label{buch:intalg:section:wurzel}}
\kopfrechts{Integrale von $R(x,\vartheta)$ mit $\vartheta=\!\sqrt{ax^2+bx+c}$}
In Kapitel~\ref{chapter:ratint} ist es uns bereits gelungen, Stammfunktionen
für beliebige rationale Funktionen zu finden.
Der Formalismus von Abschnitt~\ref{buch:intalg:section:elementar} zeigt,
dass die Stammfunktion immer ein elementare Funktion ist.
In diesem Abschnitt soll untersucht werden, was sich ändert, wenn man den
Körper $\mathbb{Q}(x)$ der rationalen Funktionen algebraisch um die
Quadratwurzel 
\[
\vartheta = \!\sqrt{ax^2+bx+c}
\]
erweitert.
Die Vorgehensweise soll später verallgemeinert werden und zeigen, wie
allgemein entschieden werden kann, ob ein Integrand eine elementare
Stammfunktion hat.

%
% Der Funktionekörper
%
\subsection{Der Funktionenkörper $\mathbb{Q}(x,\vartheta)$
\label{buch:intalg:wurzel:subsection:koerper}}
Der Körper $\mathbb{Q}(z,\vartheta)$ besteht aus den Funktionen
\[
f(x)
=
\frac{
\tilde{p}_0 + \tilde{p}_1\vartheta + \tilde{p}_2\vartheta^2
+ \ldots + \tilde{p}_n\vartheta^n
}{
\tilde{q}_0 + \tilde{q}_1\vartheta + \tilde{q}_2\vartheta^2
+ \ldots + \tilde{q}_m\vartheta^m
}
\]
mit $p_k\in\mathbb{Q}(x)$, $k=1,\dots,n$, und $q_k\in\mathbb{Q}(x)$,
$k=1,\dots,m$.
Da $\vartheta$ eine algebraisches Element über $\mathbb{Q}$ ist, welches
die Identität
\[
\vartheta^2 = ax^2 +bx +c
\]
mit $a,b,c\in\mathbb{Q}$ erfüllt, können die höheren Potenzen
$\vartheta^{2l}$ und $\vartheta^{2l+1}$ durch Potenzen des Radikanden
ersetze werden:
\begin{align*}
\vartheta^{2l}
&=
(ax^2 + bx + c)^l
&&\text{und}&
\vartheta^{2l+1}
&=
(ax^2 + bx + c)^l \vartheta.
\end{align*}
Es ist daher möglich, die Funktion $f(x)$ in der einfacheren Form
\begin{equation}
f(x)
=
\frac{p_0+p_1\vartheta}{q_0+q_1\vartheta}
\label{buch:algint:wurzel:eqn:funktion}
\end{equation}
zu schreiben.
Gesucht ist also eine Stammfunktion von
\[
f
\in
\mathbb{Q}(x,\vartheta)
=
\biggl\{
\frac{p_0+p_1\vartheta}{q_0+q_1\vartheta}
\;
\bigg|
\;
p_k,q_k\in\mathbb{Q}(x), k=0,1
\biggr\}
\]
als elementare Funktion.
Da schon die Stammfunktionen von rationalen Funktionen nicht unbedingt
rational sein müssen, kann man nicht erwarten, dass eine Stammfunktion
einer Funktion in $\mathbb{Q}(x,\vartheta)$ in diesem Körper liegen wird.
Es werden also Erweiterungen des Körpers um weitere Funktionen nötig sein,
um das Problem lösen zu können.

%
% Rationalisierung
%
\subsection{Rationalisierung
\label{buch:intalg:wurzel:subsection:rationalisierung}}
Die Darstellung~\eqref{buch:algint:wurzel:eqn:funktion} der Funktion kann
weiter vereinfacht werden, mit potenzieller Vereinfachung der Berechnung
der Stammfunktion.
Dazu wird \eqref{buch:algint:wurzel:eqn:funktion} mit
$q_0-q_1\vartheta$ zu
\begin{align}
f
&=
\frac{p_0+p_1\vartheta}{q_0+q_1\vartheta}
\cdot
\frac{q_0-q_1\vartheta}{q_0-q_1\vartheta}
\notag
\\
&=
\frac{
p_0q_0 +(p_0q_1+p_1q_0)\vartheta + p_1q_1\vartheta^2
}{
q_0^2-q_1^2\vartheta^2
}
\notag
\\
&=
\frac{
p_0q_0+p_1q_1(ax^2+bx+c)
}{
q_0^2-q_1^2(ax^2+bx+c)
}
+
\frac{
p_0q_1+p_1q_0
}{
q_0^2-q_1^2(ax^2+bx+c)
}
\vartheta
\notag
\\
&=
p + q\vartheta
\label{buch:intalg:wurzel:eqn:linkombtheta}
\end{align}
erweitert.
Daraus liest man ab, dass der Funktionenkörper sogar als
\[
\mathbb{Q}(x,\vartheta)
=
\{
p+q\vartheta
\mid
p,q\in\mathbb{Q}(x)
\}
\]
geschrieben werden kann.

Die Relation $\vartheta^2-(ax^2+bx+c)=0$ für das algebraische Element
$\vartheta$ bedeutet auch, dass
\[
\vartheta
=
\frac{ax^2+bx+c}{\vartheta}
\]
und damit, dass der Funktionenkörper auch aus den Linearkombinationen
\begin{equation}
\mathbb{Q}(x,\vartheta)
=
\biggl\{
p+q\frac{1}{\vartheta}
\bigg|
p,q\in\mathbb{Q}(x)
\biggr\}
\label{buch:intalg:wurzel:eqn:linkombthetainverse}
\end{equation}
von $1$ und $1/\vartheta$ mit Koeffizienten in $\mathbb{Q}(x)$ besteht.

%
% Das Integrationsproblem
%
\subsection{Das Integrationsproblem
\label{buch:intalg:wurzel:subsection:integration}}
Die in Abschnitt~\ref{buch:intalg:wurzel:subsection:rationalisierung}
erfolgte Rationalisierung zeigt, dass eine rationale Funktion der Form
$R(x,\vartheta)$ immer in die Form \eqref{buch:intalg:wurzel:eqn:linkombtheta}
oder \eqref{buch:intalg:wurzel:eqn:linkombthetainverse} gebracht
werden kann.
Es müssen also nur noch Integrale der Form
\begin{equation*}
\int
p(x)+q(x)\vartheta
\,dx
\qquad\text{oder}\qquad
\int
p(x)+q(x)\frac{1}{\vartheta}
\,dx
\end{equation*}
gelöst werden, wobei $p$ und $q$ rationale Funktionen in $\mathbb{Q}(x)$
sind.
Das Integral über den ersten Summanden ist das Integral einer rationalen
Funktion und kann daher mit der Methode von Kapitel~\ref{chapter:ratint}
integriert werden.
Es bleibt daher nur noch Integrale der Form
\begin{equation}
\int q(x)\,\vartheta\,dx
\qquad\text{oder}\qquad
\int \frac{q(x)}{\vartheta}\,dx
\label{buch:intalg:wurzel:eqn:integrand}
\end{equation}
zu bestimmen.


Im Folgenden bevorzugen wir die zweite Form des Integrals in
\eqref{buch:intalg:wurzel:eqn:integrand}.
Um den Grund zu verstehen, betrachten wir die Ableitung
\begin{align*}
\vartheta^2
&=
ax^2+bx+c
\\
2\vartheta\vartheta'
&=
2ax + b
\intertext{und drücken $\vartheta'$ durch $\vartheta$ als}
\vartheta'
&=
\frac{2ax+b}{2\vartheta}
=
\frac{ax+b/2}{\vartheta}
&&\Rightarrow&
\frac{1}{\vartheta}
&=
\frac{\vartheta'}{ax+b/2}
\end{align*}
aus.
Man kann also $1/\vartheta$ durch die Ableitung $\vartheta'$ ersetzen,
was die Möglichkeit der partiellen Integration eröffnet.

Die rationale Funktion $q(x)$ im rechten Integral von
\eqref{buch:intalg:wurzel:eqn:integrand} ist ein Bruch von Polynomen
von $x$.
Durch Dividieren mit Rest lässt es sich als Summe
\begin{equation}
q(x)
=
\varphi(x)
+
w(x)
\qquad
\varphi(x)\in\mathbb{Q}[x],
w(x)\in\mathbb{Q}(x)
\label{buch:intalg:wurzel:eqn:summe}
\end{equation}
eines Polynoms und einer vollständig gekürzten rationalen Funktion
schreiben, deren Nenner grösseren Grad hat als der Zähler.

%
% Integranden der Form Polynom/theta
%
\subsection{Integranden der Form $\text{Polynom}/\vartheta$}
Der Summand $\varphi$ in
\eqref{buch:intalg:wurzel:eqn:summe}
lässt sich auf ein Integral nur von $1/\vartheta$ reduzieren,
wie der folgende Satz zeigt.

\begin{satz}
\label{buch:intalg:wurzel:satz:polynomtheta}
Sei $\varphi\in\mathbb{Q}[x]$ ein Polynom mit rationalen Koeffizienten
und $\vartheta=\!\sqrt{ax^2+bx+c}$.
Dann gibt es ein Polynom $\psi\in\mathbb{Q}[x]$ mit 
\begin{equation}
\int\frac{\varphi}{\vartheta}\,dx
=
\psi\vartheta
+
A\int\frac{1}{\vartheta}\,dx,
\label{buch:intalg:wurzel:eqn:polytheta}
\end{equation}
wobei $A$ eine Konstante ist.
\end{satz}

\begin{proof}
Wenn es das Polynom $\psi$ und die Konstante $A$ gibt, dann 
muss auch die Ableitung der Gleichung~\eqref{buch:intalg:wurzel:eqn:polytheta} 
erfüllt sein, also
\begin{equation}
\frac{\varphi}{\vartheta}
=
\psi'\vartheta + \psi\vartheta'+A\frac{1}{\vartheta}.
\label{buch:intalg:wurzel:eqn:polytheta'}
\end{equation}
Die Ableitung von $\vartheta$ ist
\[
\vartheta'
=
\frac{2ax+b}{2\vartheta}
=
\frac{ax+b/2}{\vartheta}.
\]
Eingesetzt in 
\eqref{buch:intalg:wurzel:eqn:polytheta'} und mit $\vartheta$ multipliziert
erhält man die Gleichung
\begin{align*}
\varphi
&=
\psi'\vartheta^2 + \psi (ax+b/2) + A
\\
%\varphi
&=
\psi' (ax^2+bx+c) + \psi (ax+b/2) + A.
\end{align*}
Um diese Differentialgleichung für die Funktion $\psi$ zu lösen, verwenden
wir Ansätze der Form
\[
\renewcommand{\arraycolsep}{1pt}
\begin{array}{rcrcrcrcrcccrcr}
\varphi&=& \varphi_mx^m &+& \varphi_{m-1}x^{m-1} &+&   \varphi_{m-2}x^{m-2} &+&   \varphi_{m-3}x^{m-3} &+& \ldots &+& \varphi_1 x &+& \varphi_0\phantom{.}\\
\psi   &=&              & &    \psi_{m-1}x^{m-1} &+&      \psi_{m-2}x^{m-2} &+&      \psi_{m-3}x^{m-3} &+& \ldots &+&    \psi_1 x &+&    \psi_0\phantom{.}\\
\psi'  &=&              & &                      & & (m-1)\psi_{m-1}x^{m-2} &+& (m-2)\psi_{m-2}x^{m-3} &+& \ldots &+&   2\psi_2 x &+&    \psi_1.
\end{array}
\]
Durch Ausmultiplizieren und Koeffizientenvergleich entstehen die Gleichungen
\[
\renewcommand{\arraycolsep}{1pt}
\def\i#1{\rlap{$\scriptstyle #1$}\phantom{m-2}}
\begin{array}{lcrcrcrcrcrcrcr}
\varphi_{m  } &=& (m-1)a\psi_{m-1}   & &                    & &                    &+& a\psi_{m-1}   & &                   & &   \\
\varphi_{m-1} &=& (m-2)a\psi_{m-2}   &+& (m-1)b\psi_{m-1}   & &                    &+& a\psi_{m-2}   &+& (b/2)\psi_{m-1} & &   \\
\varphi_{m-2} &=& (m-3)a\psi_{m-3}   &+& (m-2)b\psi_{m-2}   &+& (m-1)c\psi_{m-1}   &+& a\psi_{m-3}   &+& (b/2)\psi_{m-2}   & &   \\
\varphi_{m-3} &=& (m-4)a\psi_{m-4}   &+& (m-3)b\psi_{m-3}   &+& (m-2)c\psi_{m-2}   &+& a\psi_{m-4}   &+& (b/2)\psi_{m-3}   & &   \\
              &\vdots &              & &                    & &                    & &               & &                   & &   \\
\varphi_{3}   &=&     2a\psi_{\i{2}} &+&     3b\psi_{\i{3}} &+&     4c\psi_{\i{4}} &+& a\psi_{\i{2}} &+& (b/2)\psi_{\i{3}} & &   \\
\varphi_{2}   &=&      a\psi_{\i{1}} &+&     2b\psi_{\i{2}} &+&     3c\psi_{\i{3}} &+& a\psi_{\i{1}} &+& (b/2)\psi_{\i{2}} & &   \\
\varphi_{1}   &=&                    & &      b\psi_{\i{1}} &+&     2c\psi_{\i{2}} &+& a\psi_{\i{0}} &+& (b/2)\psi_{\i{1}} & &   \\
\varphi_{0}   &=&                    & &                    & &      c\psi_{\i{1}} & &               &+& (b/2)\psi_{\i{0}} &+& A \\
\end{array}
\]
oder zusammengefasst
\[
\renewcommand{\arraycolsep}{1pt}
\def\i#1{\rlap{$\scriptstyle #1$}\phantom{m-2}}
\begin{array}{lcrcrcrcrcr}
\varphi_{m  } &=& \phantom{(}m\phantom{\mathstrut-0)}a\psi_{m-1}
                                     & &                             & &                    & &   \\
\varphi_{m-1} &=& (m-1)a\psi_{m-2}   &+&  ((2(m-1)+1)b/2)\psi_{m-1}  & &                    & &   \\
\varphi_{m-2} &=& (m-2)a\psi_{m-3}   &+&  ((2(m-2)+1)b/2)\psi_{m-2}  &+& (m-1)c\psi_{m-1}   & &   \\
\varphi_{m-3} &=& (m-3)a\psi_{m-4}   &+&  ((2(m-3)+1)b/2)\psi_{m-3}  &+& (m-2)c\psi_{m-2}   & &   \\
              &\vdots &              & &                             & &                    & &   \\
\varphi_{3}   &=&     3a\psi_{\i{2}} &+& 7b/2\phantom{)}\psi_{\i{3}} &+&     4c\psi_{\i{4}} & &   \\
\varphi_{2}   &=&     2a\psi_{\i{1}} &+& 5b/2\phantom{)}\psi_{\i{2}} &+&     3c\psi_{\i{3}} & &   \\
\varphi_{1}   &=&      a\psi_{\i{0}} &+& 3b/2\phantom{)}\psi_{\i{1}} &+&     2c\psi_{\i{2}} & &   \\
\varphi_{0}   &=&                    & &  b/2\phantom{)}\psi_{\i{0}} &+&      c\psi_{\i{1}} &+& A.\\
\end{array}
\]
Daraus kann man Rekursionsformeln für die Koeffizienten $\psi_k$ ableiten:
\begin{align*}
\psi_{m-1}
&=
\frac{1}{ma}\varphi_m
\\
\psi_{m-2}
&=
\frac{1}{(m-1)a}\biggl(
\varphi_{m-1}
-
\frac{(2(m-1)+1)b}{2}\psi_{m-1}
\biggr)
\\
\psi_{m-3}
&=
\frac{1}{(m-2)a}\biggl(
\varphi_{m-2}
-
\frac{(2(m-1)+1)b}{2}\psi_{m-2}
-
(m-1)c\psi_{m-1}
\biggr)
\\[-5pt]
&\phantom{a}\vdots
\\[-3pt]
\psi_{k}
&=
\frac{1}{(k+1)a}
\biggl(
\varphi_{k+1}
-
\frac{(2(k+1)+1)b}{2}\psi_{k+1}
-
(k+2)
\psi_{k+2}
\biggr)
\\[-5pt]
&\phantom{a}\vdots
\\[-3pt]
\psi_0
&=
\frac{1}{a}\biggl(
\varphi_1
-
\frac{3b}{2}\psi_1
-
2c\psi_2
\biggr)
\\
A
&=
\varphi_0 -(b/2)\psi_0 - c\psi_1.
\end{align*}
Damit sind der Reihe nach die Koeffizienten
$\psi_{m-1},\psi_{m-2},\dots,\psi_2,\psi_1,\psi_0,A$
eindeutig festgelegt, was die Existenz des Polynoms $\psi$ und der
Konstanten $A$ beweist.
\end{proof}

Der Satz besagt, dass der erste Summand in~\eqref{buch:intalg:wurzel:eqn:summe}
zu einem Integranden führt, dessen Integral sich auf das Integral von
$1/\vartheta$ reduzieren lässt.
Bei der Berechnung des zweiten Summanden treten noch weitere Arten
solcher spezieller Integrale auf.
Diese werden in Abschnitt~\ref{buch:intalg:wurzel:pol:subsection:ktheta}
auf ein Integral mit einem Integranden der Form $1/\vartheta$ zurückgeführt
und im Abschnitt~\ref{buch:intalg:wurzel:subsection:stammfunktiontheta}
berechnet.

%
% Partialbruchzeregung des gebrochenen Teils
%
\subsection{Partialbruchzerlegung des gebrochenen Teils}
Es bleibt jetzt also noch, das Integral der Form
\[
\int \frac{w(x)}{\vartheta}\,dx
\]
für eine rationale Funktion $w(x)\in\mathbb{Q}(x)$ zu berechnen.
Die rationale Funktion $w(x)$ hat eine Partialbruchzerlegung
\begin{align*}
w(z)
&=
\sum_{i=1}^n \sum_{k=1}^{n_i}\frac{A_{ki}}{(x-\xi_i)^k},
\end{align*}
wobei $\xi_i$ die Nullstellen des Nenners mit Vielfachheit $n_k$ sind.
Das Integral wird dann zu
\begin{align*}
\int \frac{w(x)}{\vartheta}\,dx
&=
\sum_{i=1}^n \sum_{k=1}^{n_i}
A_{ki}
\int
\frac{1}{(x-\xi_i)^k\vartheta}\,dx.
\end{align*}
Der gebrochene Teil kann also mit der Partialbruchzerlegung auf die
Integrale 
\begin{equation}
\int\frac{1}{\vartheta}\,dx
\qquad\text{und}\qquad
\int\frac{1}{(x-\xi)^k\vartheta}\,dx
\label{buch:intalg:wurzel:eqn:grundintegrale}
\end{equation}
zurückgeführt werden.
Zusätzlich zum Integral mit Integrand $1/\vartheta$ treten jetzt noch
weitere Integrale auf.

%
% Die Stammfunktion von 1/theta
%
\subsection{Die Stammfunktion von $1/\vartheta$
\label{buch:intalg:wurzel:subsection:stammfunktiontheta}}
Die Stammfunktion von $1/\vartheta$, das erste Integral in 
\eqref{buch:intalg:wurzel:eqn:grundintegrale},
hängt davon ab, wieviele reelle Nullstellen der Radikand hat und
ob $a$ positiv oder negativ ist.
Die Anzahl der Nullstellen ist durch das Vorzeichen der Diskriminante
$\Delta=b^2-4ac$ bestimmt.

\subsubsection{Der Fall $b^2-4ac=0$}
In diesem Fall hat der Radikand genau eine doppelte Nullstelle
bei $-b/2a$.
Für $a<0$ ist der Radikand nicht positiv, wir betrachten daher nur
den Fall $a>0$.
Die Substitution
\begin{align*}
x
&=
\frac{z}{\!\sqrt{a}}
z-\frac{b}{2a}
\qquad
\text{mit}
\qquad
dx
=
\frac{1}{\!\sqrt{a}}
dz
\\
z &= \!\sqrt{a}\biggl(x+\frac{b}{2a}\biggr)
\intertext{verwandelt den Integranden in}
ax^2+bx+c&=z^2
\intertext{und damit das Integral in}
\int\frac{1}{\vartheta}\,dx
&=
\int \frac{1}{\sqrt{z^2}} \frac{1}{\!\sqrt{a}}\,dz
=
\frac{1}{\!\sqrt{a}} \int\frac{1}{z}\,dz
=
\frac{1}{\!\sqrt{a}} \log z
\\
&=
\frac{1}{\!\sqrt{a}} \log \!\sqrt{a}\biggl(x+\frac{b}{2a}\biggr)
=
\frac{1}{\!\sqrt{a}} \log(2ax+b)
+
\frac{1}{\!\sqrt{a}} \log\frac{\sqrt{a}}{2a}
\intertext{Da die Stammfunktion nur bis auf eine Konstante bestimmt ist,
können wir den letzten Term in der Integrationskonstante absorbieren
und}
\int\frac{1}{\!\sqrt{ax^2+bx+c}}\,dx
&=
\frac{1}{\!\sqrt{a}} \log(2ax+b) + C
\end{align*}
schreiben.

\subsubsection{Der Fall $b^2-4ac>0$}
Für $b^2-4ac>0$ kann man durch die Substitution
\begin{align*}
x
&=
-\frac{b}{2a} + z\frac{\!\sqrt{b^2-4ac}}{2a}
&&\text{und}&
dx &= \frac{\!\sqrt{b^2-4ac}}{2a} \,dz
\\
z
&=
\frac{2a}{\!\sqrt{b^2-4ac}}\biggl(x+\frac{b}{2a}\biggr)
=
\frac{1}{\!\sqrt{b^2-4ac}} (2ax+b)
\end{align*}
den Radikanden in
\[
ax^2+bx+z
=
\frac{b^2-4ac}{4a}(1-z^2)
\]
umformen.
Für $a>0$ ist also
\begin{align*}
\int \frac{1}{\vartheta}\,dx
&=
\int \!\sqrt{\frac{4a}{b^2-4ac}} \frac{1}{\!\sqrt{1-z^2}}
\frac{\!\sqrt{b^2-4ac}}{2a} \,dz
\\
&=
\frac{1}{\!\sqrt{a}}
\int
\frac{1}{\!\sqrt{1-z^2}}\,dz
=
\frac{1}{\!\sqrt{a}}
\arcsin z
+C
\\
&=
\frac{1}{\!\sqrt{a}}
\arcsin
\frac{2ax+b}{\!\sqrt{b^2-4ac}}
+ C.
%
\intertext{Für $a<0$ ist der Integrand}
\vartheta
&=
\!\sqrt{\frac{b^2-4ac}{-4a}}
\!\sqrt{z^2-1}
\qquad\text{bzw.}\qquad
\!\sqrt{z^2-1}
=
\sqrt{\frac{-4a}{b^2-4ac}}\vartheta
\intertext{und damit das Integral in}
\int \frac{1}{\vartheta}\,dx
&=
\int \!\sqrt{\frac{-4a}{b^2-4ac}}\frac{1}{\!\sqrt{z^2-1}}
\frac{\!\sqrt{b^2-4ac}}{2a} \,dz
\\
&=
\frac{1}{\!\sqrt{-a}}
\int\frac{1}{\!\sqrt{z^2-1}}\,dz
=
\frac{1}{\!\sqrt{-a}}
\log(2\!\sqrt{z^2-1}+2z)
\\
&=
\frac{1}{\!\sqrt{-a}}
\log\biggl(
2\!\sqrt{\frac{-4a}{b^2-4ac}}\vartheta
+\frac{2}{\!\sqrt{b^2-4ac}}(2ax+b)
\biggr)
\\
&=
\frac{1}{\!\sqrt{-a}}
\log\bigl(
2\!\sqrt{-a}\vartheta
+
2ax+b
\bigr)
+
\frac{1}{\!\sqrt{-a}}
\log\biggl(\frac{2}{\!\sqrt{b^2-4ac}}\biggr).
\intertext{Der zweite Term ist eine Konstante, auf die es für die
Stammfunktion nicht ankommt.
Wir absorbieren ihn daher in der Integrationskonstanten und können
schreiben}
\int\frac{1}{\vartheta}\,dx
&=
\frac{1}{\!\sqrt{-a}}\log\bigl(2\sqrt{-a}\vartheta + 2ax+b\bigr)
+C.
\end{align*}

\subsubsection{Der Fall $b^2-4ac<0$}
In diesem Fall hat der Radikand keine Nullstelle.
Für $a<0$ ist der Radikand immer negativ, wir betrachten daher wieder
nur den Fall $a>0$.
Die Substitution
\begin{align*}
x&=-\frac{b}{2a}+z\frac{\!\sqrt{4ac-b^2}}{2a}
\qquad\text{und}\qquad
dx=\frac{4ac-b^2}{2a}\,dz
\\
z&=\frac{2a}{\!\sqrt{4ac-b^2}}\biggl(x+\frac{b}{2a}\biggr)
\intertext{verwandelt den Radikanden in}
ax^2+bx+c
&=
\frac{4ac-b^2}{4a}(z^2+1)
\intertext{und damit das Integral}
\int\frac{1}{\vartheta}\,dx
&=
\int
\!\sqrt{\frac{4a}{4ac-b^2}}
\frac{1}{ \!\sqrt{z^2+1}}
\frac{\!\sqrt{4ac-b^2}}{2a}
\,dz
=
\frac{1}{2\!\sqrt{a}}
\int \frac{1}{\!\sqrt{z^2+1}}\,dz
\\
&=
\frac{1}{\!\sqrt{a}}
\operatorname{arsinh} z+C
\\
&=
\frac{1}{\!\sqrt{a}}
\operatorname{arsinh} 
\frac{2a}{\!\sqrt{4ac-b^2}}\biggl(x+\frac{b}{2a}\biggr)+C
=
\frac{1}{\!\sqrt{a}}
\operatorname{arsinh} 
\frac{2ax+b}{\!\sqrt{4ac-b^2}}+C.
\end{align*}

%
% Reelle Ausdrücke für die Stammfunktion von 1/vartheta
%
\subsubsection{Reelle Ausdrücke für die Stammfunktion von $1/\vartheta$}
Wir fassen die fallweisen Rechnungen der vorangegangenen Abschnitt
ein einem Satz zusammen.

\begin{satz}
\label{buch:intalg:wurzel:satz:1theta}
Für $\vartheta = ax^2+bx+c$ mit reellen Koeffizienten $a,b,c\in\mathbb{R}$
gilt
\begin{equation}
\int\frac{1}{\vartheta}\,dx
=
\begin{cases}
\displaystyle
\frac{1}{\!\sqrt{a}}
\log(2ax+b)+C
&\text{falls $b^2-4ac=0$ und $a>0$,}
\\[10pt]
\displaystyle
\frac{1}{\!\sqrt{a}}
\arcsin\frac{2ax+b}{\!\sqrt{b^2-4ac}}+C
&\text{falls $b^2-4ac>0$ und $a>0$,}
\\[10pt]
\displaystyle
\frac{1}{\!\sqrt{-a}}
\log\bigl(2\!\sqrt{-a}\vartheta+2ax+b\bigr) + C
&\text{falls $b^2-4ac>0$ und $a<0$,}
\\[10pt]
\displaystyle
\frac{1}{\!\sqrt{a}}
\operatorname{arsinh}\frac{2ax+b}{\!\sqrt{4ac-b^2}}+C
&\text{falls $b^2-4ac<0$ und $a>0$.}
\end{cases}
\end{equation}
\end{satz}

%
% Ein gemeinsamer Ausdruck für alle Fälle
%
\subsubsection{Ein gemeinsamer Ausdruck für alle Fälle}
Die Fallunterscheidung von Satz~\ref{buch:intalg:wurzel:satz:1theta}
ist vor allem deshalb notwendig, weil wir imaginäre Werte von
$\!\sqrt{b^2-4ac}$ und $\!\sqrt{-a}$ vermeiden möchten.
Die Funktion $\arcsin$ und $\operatorname{arsinh}$ können aber
wie in Abschnitt~\ref{buch:intalg:elementar:subsection:trigo}
gezeigt auch ausschliesslich durch Logarithmen ausgedrückt werden.
Es sollte daher möglich sein, einen einzigen komplexen
Ausdruck für eine Stammfunktion von $1/\vartheta$ zu finden.

Durch Nachrechnen der Ableitung von
\[
\Theta
=
\frac{1}{\!\sqrt{a}}
\log\bigl(
2ax + b - 2\!\sqrt{a} \vartheta
\bigr)
\]
bekommt man
\begin{align*}
\Theta'
&=
-\frac{1}{\!\sqrt{a}}
\cdot
\frac{1}{2ax+b-2\!\sqrt{a}\vartheta}
\cdot
(2a - 2\!\sqrt{a}\vartheta')
\\
&=
-\frac{1}{\!\sqrt{a}}
\cdot
\frac{1}{2ax+b-2\!\sqrt{a}\vartheta}
\cdot
\biggl(
2a -2\!\sqrt{a}\frac{2ax+b}{2\vartheta}
\biggr)
\\
&=
-\frac{1}{\!\sqrt{a}}
\cdot
\frac{1}{2ax+b-2\!\sqrt{a}\vartheta}
\cdot
\frac{
4a\vartheta -2\!\sqrt{a}(2ax+b)
}{
2\vartheta
}
\\
&=
-\frac{1}{\!\sqrt{a}}
\cdot
\frac{1}{2ax+b-2\!\sqrt{a}\vartheta}
\cdot
\frac{
2\!\sqrt{a}\cdot
\bigl(
\!\sqrt{a}\vartheta
-
2ax-b
\bigr)
}{
2\vartheta
}
=
\frac{1}{\vartheta},
\end{align*}
Die Funktion $\Theta$ ist also eine Stammfunktion von $1/\vartheta$.
Sie liegt in einem Erweiterungskörper, dem die Konstante $\!\sqrt{a}$
und der Logarithmus $\log(2ax+b-2\!\sqrt{a}\vartheta)$ hinzugefügt
werden ist.
Damit der Logarithmus hinzugefügt werden kann, muss auch noch $b$
hinzugefügt werden.
Somit ist
\[
\Theta
\in
\mathbb{Q}\bigl(x,\theta,\!\sqrt{a},b,\log(\!\sqrt{a}\vartheta-2ax-b)\bigr)
\]
eine elementare Stammfunktion von $1/\vartheta$.


%
% 27-pol.tex -- Der Fall 1/(x-xi)^k\vartheta
%
% (c) 2026 Prof Dr Andreas Müller
%
\subsection{Integrale von $1/(x-\xi)^k\vartheta$
\label{buch:intalg:wurzel:pol:subsection:ktheta}}
Es muss noch gezeigt werden, dass mit den bereits entwickelten
Methoden auch das zweite Integral
in \eqref{buch:intalg:wurzel:eqn:grundintegrale},
\[
\int
\frac{1}{(x-\xi)^k\vartheta}
\,dx,
\]
behandelt werden kann.

Die Substitution
\begin{align*}
x &= \xi + \frac{1}{t} \qquad\text{mit}\qquad dx = -\frac{1}{t^2}\,dt
\intertext{verwandelt den Radikanden in}
ax^2+bx+c
&=
\frac{(a\xi^2+b\xi+c)t^2 + (2a\xi+b)\xi + a}{t^2}
\intertext{und damit den Integranden in}
\frac{1}{(x-\xi)^k\vartheta}
&=
\frac{t^k\cdot t}{
\!\sqrt{
(a\xi^2+b\xi+c)t^2 + (2a\xi+b)t + a
}
}.
\intertext{Das Integral wird damit zu}
\int\frac{1}{(x-\xi)^k\vartheta}\vartheta\,dx
&=
-
\int
\frac{t^k\cdot t}{
\!\sqrt{
(a\xi^2+b\xi+c)t^2 + (2a\xi+b)t + a
}}
\frac{1}{t^2}\,dt
\intertext{Schreibt man $A=a\xi^2+b\xi+c$, $B=2a\xi+b$ und $C=a$, wird
das Integral zu}
&=
-\int\frac{t^{k-1}}{\!\sqrt{At^2+Bt^2+C}}\,dt.
\intertext{Dies ist wieder ein Integral}
&=
-\int \frac{t^{k-1}}{\tilde{\vartheta}}\,dt
\end{align*}
der Form
$\text{Polynom}/\tilde{\vartheta}$, aber diesmal mit
\[
\tilde{\vartheta}
=
\!\sqrt{At^2+Bt+C}.
\]
Diese Art von Integral wurde in
Satz~\ref{buch:intalg:wurzel:satz:polynomtheta}
bereits auf das Integral von $1/\tilde{\vartheta}$
zurückgeführt.


%
% 28-erfahrung.tex -- Erkenntnisse
%
% (c) 2026 Prof Dr Andreas Müller
%
\subsection{Erkenntnisse}
Aus den langwierigen Rechnungen für das Integrationsproblem in
$\mathbb{Q}(x,\vartheta)$ mit $\vartheta=\!\sqrt{ax^2+bx+c}$
können wir einige Beobachtungen ableiten,
die auch für das Integrationsproblem für beliebige elementare
Funktionen gelten könnten.
\begin{enumerate}
\item 
Das Integrationsproblem spielt sich in einem Körper ab, der als
Erweiterung der rationalen Funktionen $K=\mathbb{Q}(x)$ um ein weiteres 
Element $\vartheta$ erweitert worden ist.
Der Integrand ist also eine rationale Funktion $f\in K(\vartheta)$, sie
kann als rationale Funktion 
\[
f
=
\frac{p(\vartheta)}{q(\vartheta)}
\qquad
\text{mit}
\qquad
p(\vartheta),q(\vartheta)
\in
\mathbb{Q}(x)[\vartheta]
=
K[\vartheta]
\]
geschrieben werden.
\item
Mit der Polynomdivision kann man das Integral in einen polynomiellen
und einen rationalen Teil aufteilen:
\[
\int f
=
\underbrace{
\int \text{Polynom in $\vartheta$}
}_{
\displaystyle\text{polynomieller Teil}
}
+
\underbrace{
\int\frac{\text{Rest}(\vartheta)}{q(\vartheta)}
}_{
\displaystyle\text{rationaler Teil}
}.
\]
Die Koeffizienten der Polynome sind rationale Funktionen in $K$.
\item
Die Partialbruchzerlegung kann dazu verwendet werden, den rationalen
Teil auf eine kleine Menge einfacher Grundintegrale für diese
Art von Erweiterung zu reduzieren.
Im Kapitel~\ref{chapter:ratint} war dies die Methode von Hermite.
\item
Es braucht ein spezielle Methode, um die Grundintegrale zu bestimmen.
Im Kapitel~\ref{chapter:ratint} war dies die Methode von Rothstein-Trager.
\item
Die entstehenden Stammfunktionen haben eine gemeinsame Struktur.
Sie bestehen aus einer rationalen Funktion von $\vartheta$ und
einer Linearkombination von Logarithmusfunktionen, deren Argumente
lineare Funktionen von $\vartheta$ mit Koeffizienten in $K$ sind.
\end{enumerate}
Im nächsten Abschnitt soll gezeigt, werden, dass die in Beobachtung~5
festgestellte allgemein Form der Stammfunktion tatsächlich 
immer dann vorliegt, wenn die Stammfunktion elementar ist.
Sie ist unter dem Namen Liouville-Prinzip bekannt.




